% !TEX root = ../optim-ml/OptimML.tex


%%%%%%%%%%%%%%%%%%%%%%%%%%%%%%%%%%%%%%%%%%%%%%%%%%%%%%%%%%%%%%%%%%%%%%%%
%%%%%%%%%%%%%%%%%%%%%%%%%%%%%%%%%%%%%%%%%%%%%%%%%%%%%%%%%%%%%%%%%%%%%%%%
%%%%%%%%%%%%%%%%%%%%%%%%%%%%%%%%%%%%%%%%%%%%%%%%%%%%%%%%%%%%%%%%%%%%%%%%
\section{Regularization}

When the sample size $n$ is small relative to the model dimension $p$, regularization helps control the complexity of empirical risk minimization.


%%%%%%%%%%%%%%%%%%%%%%%%%%%%%%%%%%%%%%%%%%%%%%%%%%%%%%%%%%%%%%%%%%%%%%%%
\subsection{Penalized Least Squares}

For simplicity, consider regularized least-squares regression:
\eql{\label{eq-regul-ls}
	\umin{x \in \RR^p} f_\la(x) \eqdef \frac{1}{2}\norm{Ax-y}^2 + \la R(x) 
}
where $R(x)$ is the regularizer and $\la\geq 0$ the regularization parameter. 
%
The regularizer expresses prior structure of the weight vector $x$, such as small magnitude or sparsity. The strength $\la$ can be selected by cross-validation.

We assume that $R$ is finite, nonnegative, lower semicontinuous, and coercive, i.e. $R(x) \rightarrow +\infty$ as $\norm{x} \rightarrow +\infty$.
%
The following proposition shows that as $\la\downarrow0$, regularization selects among the exact solutions when $Ax=y$ is feasible. This is particularly useful when $\ker(A)\neq\{0\}$ and there are infinitely many exact solutions.

\begin{prop}
	Assume $Ax=y$ is feasible and $\la_k>0$ tends to zero. If $x_{\la_k}$ minimizes $f_{\la_k}$, then this sequence is bounded, and every accumulation point $x^\star$ is a solution of the constrained optimization problem
	\eql{\label{eq-regul-constr}
		\umin{Ax=y} R(x).
	}
\end{prop}


\begin{proof}
	Choose an arbitrary feasible point $x_0$, so $Ax_0=y$. Optimality of $x_{\la_k}$ gives
	\eql{\label{eq-ineq-proof-regul}
			\frac{1}{2}\norm{Ax_{\la_k}-y}^2 + \la_k R(x_{\la_k}) \leq \la_k R(x_0).
	}
	Nonnegativity of the terms implies $R(x_{\la_k}) \leq  R(x_0)$. Coercivity of the regularizer makes the sequence $(x_{\la_k})_k$ bounded, since $R$ has bounded sublevel sets.
	%
	The residual also satisfies $\norm{Ax_{\la_k}-y}^2\leq2\la_k R(x_0)$, so every accumulation point obeys $Ax^\star = y$.
	%
	Along a subsequence converging to $x^\star$, lower semicontinuity gives $R(x^\star)\leq R(x_0)$. Since $x_0$ was arbitrary, $x^\star$ solves~\eqref{eq-regul-constr}.
\end{proof}

%%%%%%%%%%%%%%%%%%%%%%%%%%%%%%%%%%%%%%%%%%%%%%%%%%%%%%%%%%%%%%%%%%%%%%%%
\subsection{Ridge Regression}

Ridge regression uses the quadratic penalty $R(x)=\frac12\norm{x}_{\RR^p}^2$. For $\la>0$, the objective in~\eqref{eq-regul-ls} has a unique minimizer:
\eq{
	x_\la \eqdef \uargmin{x \in \RR^p} f_\la(x) = \frac{1}{2} \norm{Ax-y}^2_{\RR^n} +\frac\la2\norm{x}_{\RR^p}^2.
}
One has 
\eq{
	\nabla f_\la(x)=A^\top(Ax-y)+\la x,\qquad\nabla f_\la(x_\la)=0
}
Thus $x_\la$ depends linearly on $y$ and can be computed by solving a linear system. Two equivalent expressions are available.

\begin{prop}
	One has
	\begin{align}
		x_\la &= (A^\top A + \la \Id_p)^{-1} A^\top y,  \label{eq-regul-ls-1} \\
			 &= A^\top (A A^\top + \la \Id_n)^{-1} y. \label{eq-regul-ls-2} 
	\end{align}
\end{prop}

\begin{proof}
	Denoting $B \eqdef (A^\top A + \la \Id_p)^{-1} A^\top$ and $C \eqdef A^\top (A A^\top + \la \Id_n)^{-1}$, 
	one has $(A^\top A + \la \Id_p) B = A^\top$ while 
	\eq{
		(A^\top A + \la \Id_p) C = (A^\top A + \la \Id_p) A^\top (A A^\top + \la \Id_n)^{-1}
		= A^\top (A A^\top  + \la \Id_n) (A A^\top + \la \Id_n)^{-1}
		= A^\top.
	}
	Since $A^\top A + \la \Id_p$ is invertible, this gives the desired result.
\end{proof}

The systems can be solved directly, for example by Cholesky factorization, or iteratively, for example by conjugate gradients. The latter exploits the quadratic structure more effectively than ordinary gradient descent.

For $n>p$, use the smaller system~\eqref{eq-regul-ls-1}; for $n<p$, use~\eqref{eq-regul-ls-2}.

%%%%
\paragraph{Pseudo-inverse.}

As $\la\downarrow0$, $x_\la$ converges to the minimum-norm least-squares solution $A^+y$. If $Ax=y$ is feasible, this is, by~\eqref{eq-regul-constr}
\eq{
	\uargmin{Ax=y}\norm{x} .
}
If $A$ has full column rank, so $\ker(A)=\{0\}$ and $n\geq p$, $A^\top A \in \RR^{p \times p}$ is an invertible matrix, and $(A^\top A + \la \Id_p)^{-1} \rightarrow (A^\top A)^{-1}$, so that
\eq{
	x_0 = A^+ y \qwhereq A^+ \eqdef (A^\top A)^{-1}A^\top .
}
If $A$ has full row rank, so $\ker(A^\top)=\{0\}$ and $n\leq p$, the alternative formula is
\eq{
	x_0 = A^+ y \qwhereq A^+ \eqdef A^\top (AA^\top)^{-1}.
}	 
When $n=p$ and $A$ is invertible, both formulas reduce to $A^+=A^{-1}$. For any rank, the singular value decomposition (SVD) gives the Moore--Penrose pseudoinverse by inverting the nonzero singular values and leaving the zero ones unchanged.

\begin{figure}[tbp]
\centering
\includegraphics[width=\linewidth,height=.62\textheight,keepaspectratio]{chapter-optimization-and-machine-learning-advanced-topics/ml-sparsity-lq/proposed.pdf}
\caption{$\ell ^q$ balls $\enscond {x}{\sum _k |x_k|^q \leq 1}$ for varying $q$.}
\label{fig-ml-sparsity-lq}
\end{figure}


%%%%%%%%%%%%%%%%%%%%%%%%%%%%%%%%%%%%%%%%%%%%%%%%%%%%%%%%%%%%%%%%%%%%%%%%
\subsection{Lasso}

The Lasso uses an $\ell^1$ penalty:
\eq{
	R(x) = \norm{x}_1 \eqdef \sum_{k=1}^p |x_k|.
}
The penalty promotes sparsity in the solutions $x_\la$ of
\eq{
	x_\la \in \uargmin{x \in \RR^p} f_\la(x) = \frac{1}{2} \norm{Ax-y}^2_{\RR^n} + \la \norm{x}_1
}
The solutions $x_\la$ may be nonunique, and typically have many zero coefficients.
%
Figure~\ref{fig-ml-sparsity-lq} illustrates the geometry of $\ell^q$ balls. They concentrate around the coordinate axes as $q \rightarrow 0$, favoring sparse solutions, but become nonconvex for $q<1$.


Sparsity can encode a prior on the unknown, as in imaging, or select a small subset of predictive features. Feature selection can make a model easier to interpret and cheaper to evaluate.
%
Whether Lasso or ridge regression performs better depends on sparsity, correlations, noise, and the prediction task.

The objective $f_\la$ is convex, but the $\ell^1$ penalty is nonsmooth, so ordinary gradient descent does not apply at all points.
%
Section~\ref{sec-ml-ista} develops an appropriate modification of gradient descent.
%
A closed-form solution $x_\la$ is generally unavailable, but the case of an orthogonal design matrix $A$ reduces to coordinatewise thresholding.



\begin{figure}[tbp]
\centering
\BookPanel{.485}{1}{chapter-optimization-and-machine-learning-advanced-topics/ml-varspars/original/sparse-regularization-soft-objective-1.pdf}\hfill
\BookPanel{.485}{2}{chapter-optimization-and-machine-learning-advanced-topics/ml-varspars/original/sparse-regularization-soft-objective-2.pdf}
\caption{Effect of $\la$ on the penalized scalar objective $F(x)\eqdef\frac12(x-y)^2+\la|x|$.}
\label{fig-ml-varspars}
\end{figure}
\begin{figure}[tbp]
\BookContinuedFigure
\centering
\BookPanel{.485}{3}{chapter-optimization-and-machine-learning-advanced-topics/ml-varspars/original/sparse-regularization-soft-objective-3.pdf}\hfill
\BookPanel{.485}{4}{chapter-optimization-and-machine-learning-advanced-topics/ml-varspars/original/sparse-regularization-soft-objective-4.pdf}
\caption{Effect of $\la$ on the penalized scalar objective $F(x)\eqdef\frac12(x-y)^2+\la|x|$.}
\label{fig-ml-varspars-continued}
\end{figure}


\begin{prop}\label{prop-soft-tresdh}
	When $n=p$ and $A=\Id_n$, one has
	\eq{
		\uargmin{x \in \RR^p} \frac{1}{2}\norm{x-y}^2 + \la\norm{x}_1=S_\la(y)
		\qwhereq
		S_\la(x) = ( \sign(x_k) \max(|x_k|-\la,0)  )_k
	}
\end{prop}
\begin{proof}
	The separable form $f_\la(x)=\sum_k[\frac12(x_k-y_k)^2+\la|x_k|]$ reduces the problem to minimizing the scalar
	function $x \in \RR \mapsto \frac{1}{2}(x-y)^2 + \la |x|$. 
	%
	Figure~\ref{fig-ml-varspars} illustrates this minimization graphically.
	%	
	First suppose $y\geq0$. On $x<0$, $F'(x)=x-y-\la<0$, so no minimizer is negative. On $x>0$, $F'(x)=x-y+\la$ vanishes at $x=y-\la$ when $y>\la$. Otherwise, the one-sided derivatives at zero have opposite weak signs, making zero the minimizer.
	%
	The negative case follows by the symmetry $x \mapsto -x$.
\end{proof}

The nonlinear map $S_\la$ is soft thresholding.

%%%%%%%%%%%%%%%%%%%%%%%%%%%%%%%%%%%%%%%%%%%%%%%%%%%%%%%%%%%%%%%%%%%%%%%%
\subsection{Iterative Soft Thresholding}
\label{sec-ml-ista}


We derive an iterative algorithm by minimizing a quadratic surrogate.
%
We aim at minimizing
\eq{
	f(x) \eqdef \frac{1}{2} \norm{y-Ax}^2 + \la \norm{x}_1
}
For a fixed reference point $x'$, introduce the surrogate with parameter $\tau>0$:
\eq{
	f_\tau(x,x') \eqdef f(x) - \frac{1}{2}\norm{Ax-Ax'}^2 + \frac{1}{2\tau}\norm{x-x'}^2.
}
The surrogate touches the objective, $f_\tau(x,x)=f(x)$. Its additional quadratic term is
\eq{
	K(x,x') \eqdef - \frac{1}{2}\norm{Ax-Ax'}^2 + \frac{1}{2\tau}\norm{x-x'}^2=
	\frac{1}{2}\dotp{ \pa{\frac{1}{\tau}\Id_p-A^\top A} (x-x') }{x-x'}.
}
The term $K(x,x')$ is nonnegative when $\la_{\max}(A^\top A) \leq 1/\tau$, equivalently $\tau \leq 1/\norm{A}_{\text{op}}^2$, where $\norm{A}_{\text{op}} = \si_{\max}(A)$ is the operator norm.
%
Thus $f_\tau(x,x')$ is a valid surrogate:
\eq{
	f(x) \leq f_\tau(x,x'), \quad
	f_\tau(x,x)=f(x), \qandq
	f(\cdot)-f_\tau(\cdot,x') \text{ is smooth.} 
}
The surrogate is strongly convex in its first argument, so its minimizer is unique.
%
Define the iteration
\eql{\label{eq-ml-ista-surrog}
	x_{k+1} \eqdef \uargmin{x} f_{\tau}(x,x_k)
}
which ensures by construction that
\eq{
	f(x_{k+1}) \leq f(x_k).
}

\begin{prop}
	The iterates $x_k$ defined by~\eqref{eq-ml-ista-surrog} satisfy
	\eql{\label{eq-ml-ista}
		x_{k+1} = S_{\la\tau}\pa{
			x_k - \tau A^\top( A x_k - y ) 
		}
	}
	Here $S_\la(x) = (  s_\la(x_m) )_m$ applies the scalar soft-thresholding map $s_\la(r) = \sign(r) \max(|r|-\la,0)$ coordinatewise.
\end{prop}
\begin{proof}
	One has 
	\begin{align*}
		f_\tau(x,x') 
			&= \frac{1}{2}\norm{Ax-y}^2 - \frac{1}{2}\norm{Ax-Ax'}^2 + \frac{1}{2\tau}\norm{x-x'}^2 + \la \norm{x}_1 \\
			& = C + \frac{1}{2}\norm{Ax}^2-\frac{1}{2}\norm{Ax}^2+\frac{1}{2\tau}\norm{x}^2
			- \dotp{Ax}{y} + \dotp{Ax}{Ax'} - \frac{1}{\tau}\dotp{x}{x'}
			+ \la \norm{x}_1 \\
			& = C + \frac{1}{2\tau}\norm{x}^2 + \dotp{x}{ -A^\top y+A^\top A x'-\frac{1}{\tau}x' } + \la \norm{x}_1 \\
			&= C' + \frac{1}{\tau}
				\pa{
					\frac{1}{2}\norm{x-(x'-\tau A^\top(Ax'-y))}^2
					+ \tau \la \norm{x}_1.
				}
	\end{align*}
	Proposition~\ref{prop-soft-tresdh} shows that the minimizer of $f_\tau(x,x')$ is thus
	indeed $S_{\la\tau}( x' - \tau A^\top( A x' - y ) )$ as claimed.
\end{proof}

Equation~\eqref{eq-ml-ista} defines iterative soft thresholding. For $A\neq0$, the surrogate argument applies when $0<\tau\leq1/\norm{A}_{\mathrm{op}}^2$. Convergence to a Lasso solution holds for a fixed step in the larger range $0<\tau<2/\norm{A}_{\mathrm{op}}^2$. If $A=0$, every positive step is admissible.


